% GENERATED FILE. Edit canonical lessons, then run npm run generate:books.
% canonical-source-hash: fnv1a64:0092a6dba6cc124e
% canonical-lessons: MR-C235-surekh, MR-C235-bhayankar, MR-C235-ascaryakarak, MR-C235-kantalvana, MR-C235-sampurna

\chapter{Lovely, Terrible, Amazing, Boring, Whole}
\label{ch:mr-lovely-terrible-amazing-boring-whole}
\hlchaptermodality{235}

\begin{chapteropening}
\textbf{By the end of this chapter:} \emph{I can say lovely, terrible, amazing, boring and whole.}

Complete the last lesson of chapter 235: I can say lovely, terrible, amazing, boring and whole.
\end{chapteropening}

\section[\texorpdfstring{\mr{सुरेख} (surekh)}{सुरेख (surekh)}]{\mr{सुरेख} (surekh) — lovely, neat}

\label{lesson:MR-C235-surekh}



\begin{quote}
Before the new one: say the Marathi for full, then the Marathi for slow.
\end{quote}

\subsection*{You'll want to know: \mr{सुरेख}}
\textbf{\mr{सुरेख}} — \emph{surekh} — \textquotedblleft{}lovely, neat\textquotedblright{}.

\begin{cousinweb}
\textbf{\mr{सुरेख}} (\emph{surekh}) is Sanskrit \emph{su-}, \textquotedblleft{}well\textquotedblright{}, and \emph{rekhā}, \textquotedblleft{}line\textquotedblright{} (CDIAL 10810): \textquotedblleft{}well-lined\textquotedblright{}.
\end{cousinweb}

\begin{grammarlens}[title={What you've built}]
One more word for this chapter.
\end{grammarlens}

\subsection*{Your turn}
\begin{itemize}
\raggedright
  \item \emph{Say it:} \emph{surekh}
  \item \emph{Say it:} \emph{surekh}, once more
  \item \emph{Say it:} say \emph{surekh}
  \item \emph{Recall:} say the Marathi for full, then the Marathi for slow, then say \emph{surekh} again
\end{itemize}

\begin{tcolorbox}[breakable,colback=teal!4,colframe=teal!35!black,title={Before you move on}]
What does \mr{सुरेख} mean? (\textquotedblleft{}Lovely, neat\textquotedblright{}.) Say it once more.
\end{tcolorbox}
\section[\texorpdfstring{\mr{भयंकर} (bhayaṅkar)}{भयंकर (bhayaṅkar)}]{\mr{भयंकर} (bhayaṅkar) — terrible}

\label{lesson:MR-C235-bhayankar}



\begin{quote}
Before the new one: say the Marathi for slow, then the Marathi for lovely.
\end{quote}

\subsection*{You'll want to know: \mr{भयंकर}}
\textbf{\mr{भयंकर}} — \emph{bhayaṅkar} — \textquotedblleft{}terrible\textquotedblright{}.

\begin{cousinweb}
\textbf{\mr{भयंकर}} (\emph{bhayaṅkar}) is Sanskrit \emph{bhayaṅkara}, \textquotedblleft{}causing fear\textquotedblright{}, a tatsama: \emph{bhaya}, \textquotedblleft{}fear\textquotedblright{}, and \emph{kara}, \textquotedblleft{}making\textquotedblright{}.
\end{cousinweb}

\begin{grammarlens}[title={What you've built}]
One more word for this chapter.
\end{grammarlens}

\subsection*{Your turn}
\begin{itemize}
\raggedright
  \item \emph{Say it:} \emph{bhayaṅkar}
  \item \emph{Say it:} \emph{bhayaṅkar}, once more
  \item \emph{Say it:} say \emph{bhayaṅkar}
  \item \emph{Recall:} say the Marathi for slow, then the Marathi for lovely, then say \emph{bhayaṅkar} again
\end{itemize}

\begin{tcolorbox}[breakable,colback=teal!4,colframe=teal!35!black,title={Before you move on}]
What does \mr{भयंकर} mean? (\textquotedblleft{}Terrible\textquotedblright{}.) Say it once more.
\end{tcolorbox}
\section[\texorpdfstring{\mr{आश्चर्यकारक} (āścaryakārak)}{आश्चर्यकारक (āścaryakārak)}]{\mr{आश्चर्यकारक} (āścaryakārak) — amazing}

\label{lesson:MR-C235-ascaryakarak}



\begin{quote}
Before the new one: say the Marathi for lovely, then the Marathi for terrible.
\end{quote}

\subsection*{You'll want to know: \mr{आश्चर्यकारक}}
\textbf{\mr{आश्चर्यकारक}} — \emph{āścaryakārak} — \textquotedblleft{}amazing\textquotedblright{}.

\begin{cousinweb}
\textbf{\mr{आश्चर्यकारक}} (\emph{āścaryakārak}) is a compound of Sanskrit parts: \emph{āścarya}, \textquotedblleft{}wonderful\textquotedblright{} (CDIAL 1464), and \emph{kāraka}, \textquotedblleft{}making\textquotedblright{}.
\end{cousinweb}

\begin{grammarlens}[title={What you've built}]
One more word for this chapter.
\end{grammarlens}

\subsection*{Your turn}
\begin{itemize}
\raggedright
  \item \emph{Say it:} \emph{āścaryakārak}
  \item \emph{Say it:} \emph{āścaryakārak}, once more
  \item \emph{Say it:} say \emph{āścaryakārak}
  \item \emph{Recall:} say the Marathi for lovely, then the Marathi for terrible, then say \emph{āścaryakārak} again
\end{itemize}

\begin{tcolorbox}[breakable,colback=teal!4,colframe=teal!35!black,title={Before you move on}]
What does \mr{आश्चर्यकारक} mean? (\textquotedblleft{}Amazing\textquotedblright{}.) Say it once more.
\end{tcolorbox}
\section[\texorpdfstring{\mr{कंटाळवाणा} (kaṇṭāḷvāṇā)}{कंटाळवाणा (kaṇṭāḷvāṇā)}]{\mr{कंटाळवाणा} (kaṇṭāḷvāṇā) — boring}

\label{lesson:MR-C235-kantalvana}



\begin{quote}
Before the new one: say the Marathi for terrible, then the Marathi for amazing.
\end{quote}

\subsection*{You'll want to know: \mr{कंटाळवाणा}}
\textbf{\mr{कंटाळवाणा}} — \emph{kaṇṭāḷvāṇā} — \textquotedblleft{}boring\textquotedblright{}.

\begin{cousinweb}
\textbf{\mr{कंटाळवाणा}} (\emph{kaṇṭāḷvāṇā}) is \textbf{\mr{कंटाळा}}, \textquotedblleft{}boredom\textquotedblright{}, with the ending \emph{-vāṇā}, \textquotedblleft{}causing\textquotedblright{}. The origin of \textbf{\mr{कंटाळा}} is not settled.
\end{cousinweb}

\begin{grammarlens}[title={What you've built}]
One more word for this chapter.
\end{grammarlens}

\subsection*{Your turn}
\begin{itemize}
\raggedright
  \item \emph{Say it:} \emph{kaṇṭāḷvāṇā}
  \item \emph{Say it:} \emph{kaṇṭāḷvāṇā}, once more
  \item \emph{Say it:} say \emph{kaṇṭāḷvāṇā}
  \item \emph{Recall:} say the Marathi for terrible, then the Marathi for amazing, then say \emph{kaṇṭāḷvāṇā} again
\end{itemize}

\begin{tcolorbox}[breakable,colback=teal!4,colframe=teal!35!black,title={Before you move on}]
What does \mr{कंटाळवाणा} mean? (\textquotedblleft{}Boring\textquotedblright{}.) Say it once more.
\end{tcolorbox}
\section[\texorpdfstring{\mr{संपूर्ण} (sampūrṇa)}{संपूर्ण (sampūrṇa)}]{\mr{संपूर्ण} (sampūrṇa) — whole, complete}

\label{lesson:MR-C235-sampurna}



\begin{quote}
Before the new one: say the Marathi for amazing, then the Marathi for boring.
\end{quote}

\subsection*{You'll want to know: \mr{संपूर्ण}}
\textbf{\mr{संपूर्ण}} — \emph{sampūrṇa} — \textquotedblleft{}whole, complete\textquotedblright{}.

\begin{cousinweb}
\textbf{\mr{संपूर्ण}} (\emph{sampūrṇa}) is Sanskrit \emph{sampūrṇa}, \textquotedblleft{}quite full, complete\textquotedblright{}, a tatsama: \emph{sam}, \textquotedblleft{}together\textquotedblright{}, and \emph{pūrṇa}, \textquotedblleft{}full\textquotedblright{}.
\end{cousinweb}

\begin{grammarlens}[title={What you've built}]
One more word for this chapter.
\end{grammarlens}

\subsection*{Your turn}
\begin{itemize}
\raggedright
  \item \emph{Say it:} \emph{sampūrṇa}
  \item \emph{Say it:} \emph{sampūrṇa}, once more
  \item \emph{Say it:} say \emph{sampūrṇa}
  \item \emph{Recall:} say the Marathi for amazing, then the Marathi for boring, then say \emph{sampūrṇa} again
\end{itemize}

\begin{tcolorbox}[breakable,colback=teal!4,colframe=teal!35!black,title={Before you move on}]
What does \mr{संपूर्ण} mean? (\textquotedblleft{}Whole, complete\textquotedblright{}.) Say it once more.
\end{tcolorbox}
